\section*{Chapter \ref{ch:m3:crypto-data-and-infrastructure} --- Crypto Data and Infrastructure}

\begin{solution}{exo:m3:crypto-data-and-infrastructure:1}
Ignore $(95, 99)$ (already in the book); apply $(101, 103)$, the book's id becoming 103; $(105, 106)$ starts after 104, so update 104 was
missed: discard the book and resynchronise from a new snapshot.
\end{solution}

\begin{solution}{exo:m3:crypto-data-and-infrastructure:2}
When it must query historical state at arbitrary past blocks (balances, positions, contract storage then) quickly and reliably, for research,
reconciliation or tracing; a \omterm{def:m3:crypto-data-and-infrastructure:node}{full node} can regenerate some of it but slowly.
\end{solution}

\begin{solution}{exo:m3:crypto-data-and-infrastructure:3}
``10014'' ``10026'' ``9995'' ``9987'' concatenated, asks first from low to high, then bids from high to low; its CRC32 is 3\,348\,617\,501.
\end{solution}

\begin{solution}{exo:m3:crypto-data-and-infrastructure:4}
About $pK = 0.1\% \times 20 = 2\%$ of the time (1.9\% in the simulation); a typical error lasts about 20 messages (21 simulated), until the lost level
is next updated.
\end{solution}

\begin{solution}{exo:m3:crypto-data-and-infrastructure:5}
A strategy can act only on what it has received; exchange timestamps say when the venue says something happened, earlier than the firm could know it,
so decisions timed by them look into the future by the feed's latency.
\end{solution}

\begin{solution}{exo:m3:crypto-data-and-infrastructure:6}
Detect it from the block hashes, delete or mark as orphaned the records from the two abandoned blocks, re-read the new blocks at those heights, and
propagate the corrections to anything computed from them; recent records should carry their block hash and a \omterm{def:m3:blockchains-for-traders:finality}{finality} flag.
\end{solution}

\begin{solution}{exo:m3:crypto-data-and-infrastructure:7}
About 20 messages: the stale share is $pR$ and the silent-error share $pK$, equal when $R = K = 20$ (both 1.92\% in the simulation).
\end{solution}

\begin{solution}{exo:m3:crypto-data-and-infrastructure:8}
Recording is not rebuilding: the vendor's recorder can lose messages, and its dataset reflects how it handled gaps; check for crossed books, sequence
gaps and resynchronisations, and compare with a second source.
\end{solution}

\begin{solution}{pb:m3:crypto-data-and-infrastructure:1}
\textbf{1.} About 1.9\% of messages. \textbf{2.} About 21 messages. \textbf{3.} Each message updates the lost level with probability $1/20$, so the wait is
geometric with mean about 20. \textbf{4.} Quote around a wrong mid, think a level exists when it does not, or miss an arbitrage that is not there; in the
worst case trade against a phantom crossed book. \textbf{5.} A lost deletion leaves a phantom level that stays until that exact price is updated again,
which, if the price moves away, may be never. \textbf{6.} About 4.75\%. \textbf{7.} Resynchronising (50 messages) takes longer than a typical silent error
(20); the careful client trades knowledge of being stale for being stale longer. \textbf{8.} Pulls or widens its quotes and marks its data as missing.
\textbf{9.} By applying buffered events on top of the new snapshot rather than waiting for fresh ones, the stale window shrinks to the snapshot's latency.
\textbf{10.} Protection against the client's own bugs and against identifier mistakes, and verification of the book it computes. \textbf{11.} Check the update
identifiers for gaps, look for crossed or locked books, compare the top of book with trades, and compare with a second source. \textbf{12.} A lost update,
usually a missed deletion or a late snapshot, since a real venue's book cannot cross. \textbf{13.} The receive timestamp. \textbf{14.} By their exchange
timestamps corrected for each feed's measured latency, or by receive timestamps at a single location. \textbf{15.} When each resynchronisation began and ended,
so that users can drop or flag those periods. \textbf{16.} Yes: a stale book knows it cannot be trusted; a wrong one trades. \textbf{17.} A provider suffices
for occasional queries and research; latency-sensitive trading, high query volumes and independence need own nodes. \textbf{18.} Harder: many venues, each with
its own feed conventions, public clouds, no consolidated tape, chains that reorganise; easier: free, public, full-depth data from every venue and every chain.
\textbf{19.} \emph{Named result:} at a loss rate of 0.1\% and 20 levels, 1.9\% of the time silently wrong without checks (errors lasting about 21 messages),
against 4.75\% knowingly stale with resynchronisation taking 50 messages. \textbf{20.} The firm's own copy of a venue's book, as good as its handling of the
messages it did not receive.
\end{solution}

\begin{solution}{iq:m3:crypto-data-and-infrastructure:1}
Buffer the stream, fetch a snapshot, drop buffered events already contained in it, apply the rest in order checking that each event starts right after the
book's id, set levels to absolute quantities (zero removes), and resynchronise on any gap; verify with checksums where the venue publishes them.
\iqlookfor{the sequence rule and resync.}
\end{solution}

\begin{solution}{iq:m3:crypto-data-and-infrastructure:2}
The exchange timestamp is the venue's time of the event (or message); the receive timestamp is when the firm got it. Use receive times for what a strategy knew
and exchange times to align venues and measure latency.
\iqlookfor{knowledge versus alignment.}
\end{solution}

\begin{solution}{iq:m3:crypto-data-and-infrastructure:3}
Read blocks from own nodes; decode swap logs by pool address and signature; store with block number, hash, log index and \omterm{def:m3:blockchains-for-traders:finality}{finality}; handle reorganisations by
comparing parent hashes and rewriting orphaned ranges; backfill from an \omterm{def:m3:crypto-data-and-infrastructure:node}{archive node}; version decoders with contract upgrades.
\iqlookfor{reorg handling and decoding discipline.}
\end{solution}

\begin{solution}{iq:m3:crypto-data-and-infrastructure:4}
Log the book and the message at each failure; check the formatting rule (decimal places, trailing zeros, level count), ordering of levels, the handling of zero
quantities and of levels beyond the snapshot's depth, and message loss or reordering around the failures; reproduce offline from the recorded stream.
\iqlookfor{formatting first, then sequence.}
\end{solution}

\begin{solution}{iq:m3:crypto-data-and-infrastructure:5}
Compare exchange and receive timestamps on quiet messages, measure round trips of cheap requests from machines in several regions, and watch the distribution rather
than the mean, especially in busy periods; synchronise the firm's clocks first.
\iqlookfor{clock discipline and distributions.}
\end{solution}

\begin{solution}{iq:m3:crypto-data-and-infrastructure:6}
Per-venue adapters that normalise messages and keep both timestamps; book builders with sequence and checksum checks; per-chain node clusters and \omterm{def:m3:crypto-data-and-infrastructure:log}{indexers}; a
normalised bus to strategies with staleness flags; recording of raw messages for replay; monitoring of gaps, latency and resyncs; and deployment in the regions
nearest each venue.
\iqlookfor{normalise, verify, flag and record.}
\end{solution}
